\documentclass[11pt]{article}

\usepackage[T1]{fontenc}
\usepackage[utf8]{inputenc}
\usepackage{lmodern}
\usepackage{microtype}
\usepackage[margin=1in]{geometry}
\usepackage{amsmath,amssymb,amsthm,mathtools}
\usepackage{enumitem}
\usepackage[hidelinks]{hyperref}
\usepackage[nameinlink,capitalise]{cleveref}

\newcommand{\R}{\mathbb{R}}
\newcommand{\C}{\mathbb{C}}
\newcommand{\cA}{\mathcal{A}}
\newcommand{\cH}{\mathcal{H}}
\newcommand{\cI}{\mathcal{I}}
\newcommand{\cJ}{\mathcal{J}}
\newcommand{\cD}{\mathcal{D}}
\newcommand{\id}{\mathrm{id}}
\newcommand{\MTT}{\mathrm{MTT}}

\newtheorem{theorem}{Theorem}[section]
\newtheorem{proposition}[theorem]{Proposition}
\newtheorem{lemma}[theorem]{Lemma}
\newtheorem{definition}[theorem]{Definition}
\newtheorem{assumption}[theorem]{Assumption}
\newtheorem{corollary}[theorem]{Corollary}
\theoremstyle{remark}
\newtheorem{remark}[theorem]{Remark}

\title{\Large\bf
Modal Fixed Points, Bell's Beables, and the Limits of Factorization:\\[0.25em]
Upper-Local Dynamics, Nonseparable States, and Operational No-Signaling in MTT}
\author{Peter Nero}
\date{July 2026}

\begin{document}
\maketitle

\begin{abstract}
Bell experiments do not require a controllable signal between spacelike
separated laboratories, but they do exclude a complete measurement-independent
hidden-variable description whose conditional probabilities factorize.  This
paper formulates that distinction for Modal Triplet Theory (MTT).  The upper
description is a local algebra or field theory on
$M=Y^4\times X^6$, while an admissible fixed point may define a globally
nonseparable state.  Under an explicit base-local descent hypothesis,
microcausality passes to the four-dimensional observable algebras.  Under
standard local completely positive instruments, this gives operational
no-signaling.  If the descended statistics violate CHSH while measurement
independence is retained, Bell factorization must nevertheless fail.  The
result is therefore not a Bell-local hidden-variable completion.  It is a
conditional account in which upper-local dynamics and globally constrained
state selection coexist, so the correlations need not be produced by a
spacelike dynamical influence.  A singlet packet is included as a benchmark,
not as an MTT prediction.  A physical MTT realization still requires a selected
fixed-point state, a locality-preserving descent, local instruments, and a
probability source that yields the benchmark statistics.
\end{abstract}

\section*{Revision note}

\paragraph{Supersedes.}
Version 2 supersedes the 2025 version associated with Zenodo record 17076301.

\paragraph{Reason.}
The earlier text treated projection as creating Bell correlations from an
otherwise Bell-local completion and did not keep microcausality, operational
no-signaling, state nonseparability, measurement independence, and Bell
factorization logically separate.

\paragraph{Resolution.}
This version distinguishes those notions, removes the incorrect inference
that an ordinary Kaluza--Klein reduction forces state factorization, and
reclassifies the singlet and quantum-network calculations as conditional
benchmarks until their probability and state data are selected by MTT.

\paragraph{Retained content.}
The fixed-point perspective and the proposal that apparently nonlocal
correlations may descend from an upper description with local dynamics are
retained, now with their exact assumptions and completion obligations stated.

\paragraph{Remaining boundary.}
MTT has not yet selected from first principles the physical preparation
state, the Born probability functional, and the detector instruments needed
to derive the singlet correlations rather than use them as a benchmark.

\tableofcontents

\section{Introduction: the central picture}

The familiar phrase ``spooky action at a distance'' combines two very
different claims:

\begin{enumerate}[label=(\roman*)]
\item a choice in one laboratory sends a controllable physical influence to a
spacelike separated laboratory;
\item the joint probabilities cannot be screened off into independent local
responses by conditioning on a complete hidden state.
\end{enumerate}

Bell's theorem concerns the second claim.  Relativistic no-signaling concerns
the first.  Quantum theory can violate a Bell inequality while forbidding the
use of the correlations as a faster-than-light communication channel
\cite{Bell1964,CHSH1969}.  Algebraic quantum field theory sharpens the same
distinction: spacelike observable algebras commute, yet a state on their joint
algebra need not factorize \cite{Haag1992,SummersWerner1987}.

The proposed MTT perspective is consequently not that Bell's theorem has been
evaded.  It is that the upper description can remain local at the level of
equations and observable algebras while its selected admissible state is
global and nonseparable.  Projection then reveals correlations already carried
by that state; it does not manufacture them from a separable Bell-local state.
In plain language, the two laboratories read different parts of one globally
constrained preparation; neither measurement has to send a message to the
other laboratory in order for their records to be correlated.

\subsection{What is proved and what is conditional}

This paper proves four statements at the declared level:

\begin{enumerate}[label=(\arabic*)]
\item a base-local internal descent preserves spacelike commutativity;
\item local nonselective instruments imply operational no-signaling;
\item measurement independence plus Bell factorization implies the CHSH
bound;
\item therefore any MTT packet reproducing a CHSH violation while preserving
measurement independence must be globally nonfactorizing.
\end{enumerate}

The paper does \emph{not} derive a particular entangled state, Born
probabilities, the value $2\sqrt{2}$, or a physical detector instrument from
fixed-point language alone.  Those are source rows in the completion contract
of \cref{sec:completion}.

\section{Three meanings of locality}
\label{sec:localities}

\begin{definition}[Upper algebraic locality]
Let $Y^4$ be a supplied causal spacetime and $X^6$ a compact internal
space.  Set $M=Y^4\times X^6$.  For each causally suitable
$\mathcal O\subset Y^4$, let
\[
  \cA_M(\widetilde{\mathcal O}),
  \qquad
  \widetilde{\mathcal O}:=\mathcal O\times X^6,
\]
be an upper observable algebra.  Upper algebraic locality means
\[
 [A,B]=0
 \quad\text{whenever}\quad
 A\in\cA_M(\widetilde{\mathcal O}_A),\
 B\in\cA_M(\widetilde{\mathcal O}_B),
 \quad
 \mathcal O_A\perp\mathcal O_B .
\]
\end{definition}

The product notation is a control model.  It does not assert that the internal
space is literally three independent manifolds or that its rank filtration
adds new spacetime coordinates.  In the current MTT geometric program,
$X^6$ must be supplied by the selected physical branch, and the $1<2<3$
labels are operator or lane data on that geometry.

\begin{definition}[Operational no-signaling]
For settings $a,b$ and outcomes $\alpha,\beta$, a probability packet
$p(\alpha,\beta\mid a,b)$ is no-signaling when
\[
 \sum_\beta p(\alpha,\beta\mid a,b)
 \quad\text{is independent of }b,
\qquad
 \sum_\alpha p(\alpha,\beta\mid a,b)
 \quad\text{is independent of }a.
\]
\end{definition}

\begin{definition}[Bell factorization]
Let $\lambda$ denote a proposed complete ontic state.  Bell factorization is
the screening-off condition
\[
 p(\alpha,\beta\mid a,b,\lambda)
 =
 p_A(\alpha\mid a,\lambda)\,
 p_B(\beta\mid b,\lambda)
\]
for almost every $\lambda$.  Measurement independence is
\[
 \rho(\lambda\mid a,b)=\rho(\lambda).
\]
\end{definition}

These notions are not interchangeable.  Commuting algebras need not carry a
product state.  No-signaling is an observed marginal condition, whereas Bell
factorization is a conditional-independence claim about a complete underlying
description.  A theory may satisfy the first two notions and fail the third.

\section{The upper fixed-point packet}
\label{sec:upper-packet}

The minimal MTT packet needed in this paper consists of:

\begin{enumerate}[label=(U\arabic*)]
\item a causal base $Y^4$, an internal space $X^6$, and an upper local
net $\mathcal O\mapsto\cA_M(\mathcal O\times X^6)$;
\item preparation data $\lambda$ with a measure $\rho$ selected before the
laboratory settings;
\item an admissible fixed point or stationary upper state
$\omega_\lambda$ selected from the preparation data, not from the later
choices $a,b$;
\item a base-local descent $\cD$ from upper observables to the effective
four-dimensional local algebras;
\item local instruments $\{\cI_{a}^{\alpha}\}_{\alpha}$ and
$\{\cJ_{b}^{\beta}\}_{\beta}$;
\item a probability rule evaluating the descended state on those
instruments.
\end{enumerate}

The phrase ``fixed point'' may refer to a fixed point of an auxiliary
stabilization flow.  Its parameter is not automatically physical time.
Likewise, existence of a fixed point does not by itself imply that the
associated Lorentzian theory is causal or that its state is entangled.  The
local-net and state hypotheses are separate rows.

\subsection{Why the state must not depend silently on both settings}

An expression such as
\[
 \Psi^\ast=\Psi^\ast(a,b,\xi)
\]
cannot simultaneously be treated as a complete premeasurement hidden state
and as independent of $a,b$.  There are only four coherent readings:

\begin{enumerate}[label=(\alph*)]
\item $\xi$ is incomplete and the complete state includes
$\Psi^\ast(a,b,\xi)$;
\item the settings are boundary data in an atemporal global boundary-value
problem;
\item the construction is retrocausal;
\item measurement independence is abandoned.
\end{enumerate}

This paper uses a fifth, cleaner preparation branch: a common preparation
selects $\omega_\lambda$ before $a,b$, and the settings choose only local
instruments.  The state may be nonseparable, but it is not recalculated from
the two later settings.  This is compatible with measurement independence and
makes the remaining failure of Bell factorization explicit.

\section{Locality-preserving descent}
\label{sec:descent}

Let $\Pi_x$ be an internal coherent projector acting fiberwise over
$x\in Y^4$.  Schematically, a descended observable has the form
\[
 \cD_{\mathcal O}(A)
 =
 \int_{X^6}(\Pi A\Pi)(\,\cdot\,,u)\,d\nu_X(u),
\]
where the integral is a bounded weak operator integral.  The crucial
condition is not merely boundedness of $\Pi$: $\Pi$ and the integral must
preserve support in the base region $\mathcal O$.

\begin{assumption}[Base-local descent]
\label{ass:base-local}
For every $\mathcal O\subset Y^4$,
\[
 \cD_{\mathcal O}
 \bigl(\cA_M(\mathcal O\times X^6)\bigr)
 \subseteq \cA_Y(\mathcal O).
\]
For spacelike $\mathcal O_A,\mathcal O_B$, representatives of the two weak
operator integrals commute pairwise at every pair of internal points.
\end{assumption}

\begin{lemma}[Locality descent]
\label{lem:locality-descent}
Under upper algebraic locality and
\cref{ass:base-local},
\[
 [\cD_{\mathcal O_A}(A),\cD_{\mathcal O_B}(B)]=0
\]
for $A\in\cA_M(\widetilde{\mathcal O}_A)$ and
$B\in\cA_M(\widetilde{\mathcal O}_B)$ whenever
$\mathcal O_A\perp\mathcal O_B$.
\end{lemma}

\begin{proof}
Using bounded weak integration and bilinearity of the commutator,
\[
\begin{split}
 [\cD_{\mathcal O_A}(A),\cD_{\mathcal O_B}(B)]
 &=
 \int_{X^6}\!\!\int_{X^6}
 [(\Pi A\Pi)(u),(\Pi B\Pi)(v)]\,
 d\nu_X(u)d\nu_X(v).
\end{split}
\]
Every integrand vanishes by upper locality and base-support preservation.
The integral is therefore zero.
\end{proof}

\begin{remark}
A generic spatially nonlocal projector does not satisfy
\cref{ass:base-local}.  Nor does a global fixed-point constraint automatically
prove it.  This is why locality descent must be stated as a theorem hypothesis
and verified for the selected MTT operator.
\end{remark}

\section{Local instruments and no-signaling}
\label{sec:instruments}

Let Alice's and Bob's local instruments be normal completely positive maps.
Write their Heisenberg dual operations as
$\cI_{a}^{\alpha\ast}$ and $\cJ_{b}^{\beta\ast}$.  The corresponding
nonselective operations are
\[
 \cI_a^\ast=\sum_\alpha\cI_{a}^{\alpha\ast},
 \qquad
 \cJ_b^\ast=\sum_\beta\cJ_{b}^{\beta\ast}.
\]
They are unital.  Spacelike locality means that Bob's nonselective operation
acts trivially on Alice's local observables, and conversely
\cite{DaviesLewis1970,Haag1992}.

\begin{proposition}[Operational no-signaling from local instruments]
\label{prop:no-signaling}
Let $\omega_\lambda$ be any normal, possibly nonseparable state on the joint
algebra.  If the two instruments are local in the preceding sense, then
\[
 p_\lambda(\alpha,\beta\mid a,b)
 :=
 \omega_\lambda\!\left(
 \cI_{a}^{\alpha\ast}
 \cJ_{b}^{\beta\ast}(\mathbf 1)\right)
\]
is operationally no-signaling.  The same remains true after averaging over a
setting-independent measure $\rho(d\lambda)$.
\end{proposition}

\begin{proof}
Summing over Bob's outcomes gives
\[
\begin{split}
\sum_\beta p_\lambda(\alpha,\beta\mid a,b)
 &=
\omega_\lambda\!\left(
\cI_{a}^{\alpha\ast}\cJ_b^\ast(\mathbf 1)\right)\\
 &=
\omega_\lambda\!\left(
\cI_{a}^{\alpha\ast}(\mathbf 1)\right),
\end{split}
\]
which is independent of $b$.  The other marginal is analogous.  Integration
against $\rho(d\lambda)$ preserves both equalities.
\end{proof}

This result locates the absence of action at a distance precisely.  A remote
choice cannot control the local marginal.  The proposition does not state
that the joint state is separable or that its conditional probabilities
factorize.

\section{Bell's boundary}
\label{sec:bell}

For dichotomic outcomes in $\{-1,+1\}$, define
\[
 E(a,b)
 =
 \sum_{\alpha,\beta=\pm1}
 \alpha\beta\,p(\alpha,\beta\mid a,b)
\]
and
\[
 S=E(a,b)+E(a,b')+E(a',b)-E(a',b').
\]

\begin{theorem}[CHSH under measurement independence and factorization]
\label{thm:chsh}
Suppose $\rho(d\lambda)$ is independent of $a,b$ and the conditional
probabilities factorize for almost every $\lambda$.  Then
\[
 |S|\leq 2.
\]
\end{theorem}

\begin{proof}
Set
\[
 A_a(\lambda)=\sum_{\alpha=\pm1}\alpha\,
 p_A(\alpha\mid a,\lambda),
\qquad
 B_b(\lambda)=\sum_{\beta=\pm1}\beta\,
 p_B(\beta\mid b,\lambda),
\]
so $|A_a|,|B_b|\leq1$.  Factorization gives
$E(a,b)=\int A_a(\lambda)B_b(\lambda)\rho(d\lambda)$.  For every
$\lambda$,
\[
\begin{split}
 &\left|
 A_a(B_b+B_{b'})
 +A_{a'}(B_b-B_{b'})
 \right|\\
 &\hspace{4em}\leq
 |B_b+B_{b'}|+|B_b-B_{b'}|
 \leq2.
\end{split}
\]
Integration proves the result.
\end{proof}

\begin{corollary}[Required MTT nonfactorization]
\label{cor:required-nonfactorization}
If a selected MTT packet preserves measurement independence and produces
$|S|>2$, Bell factorization fails on a set of nonzero
$\rho$-measure.
\end{corollary}

\begin{proof}
Otherwise \cref{thm:chsh} would imply $|S|\leq2$.
\end{proof}

This is not a defect to be hidden.  It is the mathematical content of the
proposal: MTT may retain upper algebraic locality and no-signaling while its
admissible state fails the classical screening-off condition.  In Jarrett's
language, one must not conflate observed parameter independence with the full
factorization package \cite{Jarrett1984,Fine1982}.

\section{The singlet benchmark}
\label{sec:benchmark}

Take
\[
 \cH=\C^2\otimes\C^2,
\qquad
 |\psi^-\rangle
 =
 \frac{|01\rangle-|10\rangle}{\sqrt2}.
\]
Choose
\[
 A=\sigma_z,\qquad A'=\sigma_x,
\]
\[
 B=\frac{\sigma_z+\sigma_x}{\sqrt2},
\qquad
 B'=\frac{\sigma_z-\sigma_x}{\sqrt2}.
\]
The standard quantum packet gives
\[
 |S|=2\sqrt2
\]
while its local marginals are independent of the remote settings
\cite{Tsirelson1980}.  The observables on opposite tensor factors commute.
This is the standard example of local commutativity, no-signaling, and state
nonfactorization coexisting.

\begin{proposition}[Conditional MTT realization]
\label{prop:conditional-realization}
Suppose one selected MTT source emits:

\begin{enumerate}[label=(\alph*)]
\item a preparation-selected upper fixed-point state whose descent is
$|\psi^-\rangle\langle\psi^-|$;
\item the four displayed local measurement instruments;
\item the probability functional used in
\cref{prop:no-signaling};
\item the locality hypotheses of \cref{lem:locality-descent}.
\end{enumerate}

Then the descended packet is no-signaling, reaches
$|S|=2\sqrt2$, and cannot satisfy Bell factorization under measurement
independence.
\end{proposition}

\begin{proof}
No-signaling follows from \cref{prop:no-signaling}; the CHSH value is the
benchmark calculation above; nonfactorization follows from
\cref{cor:required-nonfactorization}.
\end{proof}

The proposition is a commuting completion contract, not a derivation of its
premises.  In particular, choosing MTT data because their descent equals the
singlet would be a reconstruction or replay, not an independent prediction.

\section{Interpretation and consequences: what projection explains}

Projection can make the origin of correlations less visible.  Two distant
four-dimensional observables may be restrictions of one globally constrained
upper state.  From the lower description alone the correlations can then look
like a direct relation between the outcomes, even though no signal is sent
during the measurements.

Projection alone, however, does not select:

\begin{itemize}
\item a nonseparable upper state;
\item a probability measure on preparations;
\item the Born functional;
\item the local instrument family;
\item the strength or sign of a Bell correlator.
\end{itemize}

The accurate phrase is therefore \emph{descent of upper nonseparability}, not
``creation of nonlocality by projection.''  The explanatory work belongs to
the source theorem that selects the global admissible state and its
probability/instrument packet.

\section{Comparison with Kaluza--Klein and local QFT}

An ordinary Kaluza--Klein reduction can preserve base locality, but locality
does not force the state to factorize.  A local four-dimensional quantum field
theory may possess entangled states and Bell-violating correlations while its
spacelike algebras commute.  The earlier claim that zero-mode reduction by
itself yields Bell factorization was therefore incorrect.

The possible MTT contribution is different.  MTT proposes that a restricted
class of globally coherent states is selected by admissibility and fixed-point
conditions.  If that selection can be derived from the same upper geometry
that supplies the local net and descent, it could explain why a particular
nonseparable state occurs.  Until that source theorem is supplied, MTT and
standard local quantum theory share the same operational reconciliation of
Bell violation with no-signaling; MTT adds a conditional selection mechanism,
not yet a completed prediction.

\section{Networks and multipartite experiments}

The distinction scales to several laboratories.  Pairwise commuting local
algebras can carry a globally nonseparable state, and local instruments then
preserve operational no-signaling.  Classical network inequalities impose
additional source-independence factorizations.  A violation shows that at
least one declared network factorization fails; it does not prove a
superluminal signal.

No theorem in this paper states that MTT realizes every quantum network
correlation.  Such a claim would require a selected multipartite state,
source-independence conventions, local instruments, and the corresponding
probability functional for each network.

\section{Relation to temporal Bell inequalities}

Spatial Bell and temporal Leggett--Garg experiments test different contracts.
The spatial case uses separated laboratories, two-way no-signaling, and Bell
factorization.  The temporal case uses sequential instruments on one system
and permits forward disturbance.  The revised temporal Bell paper therefore
treats its Leggett--Garg packet separately \cite{MTTTemporalBell}.

A common upper fixed-point language may eventually connect the two, but a
unification theorem would have to preserve both operational probability laws
and their different causal constraints.  Merely using the word ``projection''
in both settings does not provide that theorem.

\section{Completion and falsification contract}
\label{sec:completion}

For the upper-local account to become a physical MTT result, one same-source
packet must provide:

\begin{enumerate}[label=(C\arabic*)]
\item a selected upper Lorentzian local net or hyperbolic field system;
\item a preparation-selected admissible fixed-point state;
\item proof that the preparation measure is independent of later settings;
\item a base-local coherent projector and descent satisfying
\cref{ass:base-local};
\item local completely positive instruments representing the actual
detectors;
\item a selected probability functional;
\item a calculated CHSH packet with uncertainty and provenance.
\end{enumerate}

The proposal fails in the following clear cases:

\begin{itemize}
\item the selected descent spreads support across spacelike base regions;
\item the remote setting changes an observed local marginal;
\item the complete state is claimed to be measurement-independent and
factorizing while $|S|>2$;
\item the state or probability rows are inserted from observed Bell data and
then reported as predictions;
\item the fixed-point solution depends on both settings without declaring a
global-boundary, retrocausal, contextual, or measurement-dependent branch.
\end{itemize}

\section{Conclusion}

The upper-world perspective can remove the need for a dynamical ``spooky
action'' between the measurements, but it does not restore Bell
factorization.  The coherent statement is:
\[
\boxed{
\text{upper-local dynamics}
+
\text{microcausal descent}
+
\text{local instruments}
+
\text{global nonseparable state}.
}
\]
This package yields no operational superluminal signal.  If it reproduces a
Bell violation while measurement independence is retained, its complete
probabilities are necessarily nonfactorizing.

That is already a useful reconciliation.  What remains specifically MTT is
to derive the nonseparable fixed-point state, the probability functional, and
the detector instruments from one selected upper source rather than importing
the quantum benchmark.  The present paper states that target without
mistaking a perspective shift for a completed source theorem.

\bibliographystyle{plain}
\bibliography{main}

\end{document}
